% ============================================================
%  Section 9: Conclusion
% ============================================================

\section{Conclusion}

% ---- Key takeaways ----------------------------------------
\begin{frame}{Key Takeaways}
  \begin{enumerate}
    \item \textbf{Right tool for the task}: Inline completion, chat, and agentic tools serve different needs. Match capability to workflow rather than defaulting to one tool for everything.

    \vspace{0.4em}
    \item \textbf{Agentic oversight is mandatory}: Agents are powerful but unreliable on consequential, irreversible actions. Build in human checkpoints.

    \vspace{0.4em}
    \item \textbf{Literature tools amplify, don't replace}: Elicit, ResearchRabbit, and Semantic Scholar accelerate discovery; human judgment determines relevance and quality.

    \vspace{0.4em}
    \item \textbf{You remain the author}: AI writing assistance is legitimate; AI authorship is not. Use tools to improve your writing, not to replace your thinking.

    \vspace{0.4em}
    \item \textbf{Hallucination vigilance}: Verify every citation, every factual claim, every code snippet. Fluency is not correctness.
  \end{enumerate}
\end{frame}

% ---- Researcher's heuristic --------------------------------
\begin{frame}{A Researcher's Heuristic}
  \centering
  \begin{tikzpicture}[scale=0.95]
    % Axes
    \draw[-{Stealth[length=6pt]}, thick] (-0.3,0) -- (8.3,0)
      node[right, font=\small] {Task complexity $\rightarrow$};
    \draw[-{Stealth[length=6pt]}, thick] (0,-0.3) -- (0,5.8)
      node[above, font=\small] {Data sensitivity $\uparrow$};

    % Quadrant boxes
    \fill[blue!10]   (0.1,0.1) rectangle (4.0,2.8);
    \fill[green!10]  (4.1,0.1) rectangle (8.0,2.8);
    \fill[orange!15] (0.1,2.9) rectangle (4.0,5.6);
    \fill[red!10]    (4.1,2.9) rectangle (8.0,5.6);

    % Quadrant labels
    \node[align=center, font=\small] at (2.0, 1.6)
      {\textbf{Low sensitivity,}\\Low complexity:\\\tool{Claude} / \tool{Codex}\\Chat};
    \node[align=center, font=\small] at (6.0, 1.6)
      {\textbf{Low sensitivity,}\\High complexity:\\\tool{Claude Code} /\\Agentic workflow};
    \node[align=center, font=\small] at (2.0, 4.2)
      {\textbf{High sensitivity,}\\Low complexity:\\Local \tool{Ollama}\\(small model)};
    \node[align=center, font=\small] at (6.0, 4.2)
      {\textbf{High sensitivity,}\\High complexity:\\Local \tool{vLLM}\\(70B+ model)};

    % Midlines
    \draw[dashed, gray] (4.0, 0) -- (4.0, 5.8);
    \draw[dashed, gray] (0, 2.8) -- (8.0, 2.8);
  \end{tikzpicture}
\end{frame}

% ---- Resources --------------------------------------------
\begin{frame}{Resources \& Further Reading}
  \begin{columns}[T, onlytextwidth]
    \begin{column}{0.55\textwidth}
      \textbf{Model documentation}\\[0.2em]
      \begin{itemize}
        \item Anthropic Claude docs: \href{https://docs.anthropic.com}{docs.anthropic.com}
        \item OpenAI: \href{https://platform.openai.com/docs}{platform.openai.com/docs}
        \item Hugging Face model hub: \href{https://huggingface.co/models}{huggingface.co/models}
      \end{itemize}
      \vspace{0.5em}
      \textbf{Key papers}\\[0.2em]
      \begin{itemize}
        \item ReAct: \parencite{yao2023react}
        \item Chain-of-Thought: \parencite{wei2022cot}
        \item SWE-Bench: \parencite{jimenez2024swebench}
        \item Foundation Model risks: \parencite{bommasani2021foundation}
        \item Hallucination survey: \parencite{huang2023hallucination}
      \end{itemize}
    \end{column}
    \hfill
    \begin{column}{0.42\textwidth}
      \textbf{Tools covered today}\\[0.2em]
      \begin{itemize}
        \item MCP spec: \href{https://modelcontextprotocol.io}{modelcontextprotocol.io}
        \item Claude Code: \href{https://claude.ai/code}{claude.ai/code}
        \item Elicit: \href{https://elicit.com}{elicit.com}
        \item ResearchRabbit: \href{https://www.researchrabbit.ai}{researchrabbit.ai}
        \item Semantic Scholar API: \href{https://api.semanticscholar.org}{api.semanticscholar.org}
        \item Ollama: \href{https://ollama.ai}{ollama.ai}
      \end{itemize}
      \vspace{0.4em}
      \centering
      \qrcode[height=1.2cm]{https://github.com/your-repo/ai-tools-for-research}\\
      \tiny Slides repository
    \end{column}
  \end{columns}
\end{frame}

% ---- Call to action ----------------------------------------
\begin{frame}{Three Concrete Next Steps}
  \begin{enumerate}
    \item \textbf{Today:} Install \tool{Claude Code} or \tool{VS-Codex} and use it for your next coding task. Pay attention to where it helps and where it errs.

    \vspace{0.6em}
    \item \textbf{This week:} Set up a \tool{ResearchRabbit} collection with 3 seed papers from your current project. Explore what it surfaces.

    \vspace{0.6em}
    \item \textbf{This month:} For your next paper submission, draft one section with AI assistance. Document your process — what prompts you used, what you changed, what you verified.
  \end{enumerate}

  \vspace{0.8em}
  \begin{center}
    \small\textit{The goal is not to use AI everywhere. It is to use it \textbf{deliberately} where it genuinely accelerates your research.}
  \end{center}
\end{frame}

% ---- Questions slide ----------------------------------------
\begin{frame}[standout]
  \Large Questions \& Discussion

  \vspace{1.5em}
  \normalsize
  David Adams\\
  \href{mailto:david.adams1@student.unimelb.edu.au}{david.adams1@student.unimelb.edu.au}

  \vspace{1.2em}
  \qrcode[height=1.5cm]{https://github.com/dadams2/ai_tools_for_research}\\[0.4em]
  \small Slides available at the link above
\end{frame}
